% =====================================================================
\section{Results, Limitations and Conclusion}\label{sec:results}
\subsection{What was measured}
\begin{table}[h]
\centering\small
\caption{Quantities taken from the program output of this report's build (not from memory).}
\label{tab:measure}
\begin{tabular}{@{}p{8.7cm}p{7.2cm}@{}}
\toprule
Quantity & Value (source)\\\midrule
Leaf parts in the scene, in 14 export objects & \nLeaves{} (\code{scene\_export.json}, Table~\ref{tab:parts})\\
Parts with emissive $>0$ & \nEmissive{} (\code{scene\_export.json})\\
Track links of the tank & \nTrackLinks{} (export; T80)\\
Mesh memory, indexed vs non-indexed & \memIndexed{} B vs \memNonIndexed{} B (Table~\ref{tab:memory})\\
Automated tests & \nTests{} run, \nPass{} pass, \nFail{} fail (\code{test\_results.json})\\
Terrain & 49\,921 vertices, 97\,440 triangles, 36\,m radius (T13, T61)\\
Bunker clear height / walk eye height & 2.60\,m / 1.7\,m (T62, T51)\\
Missile: first launch request after lock & 1.53\,s (T108); detonation 0.50\,m from the spot (T100)\\
Shell range on level ground & 203.8\,m, $-0.24$\,\% from theory (T101)\\
Debris: pieces at rest within 8\,s & 23 of 23 (T106)\\
180\,s headless run: UAV max radius, min clearance, rover hidden & 15.57\,m, 5.62\,m, 61\,\% of frames (T45)\\
Full base, 120\,s simulated & 9 controllables (7 possessable), finite state (T113)\\
\textbf{FPS / ms per frame on the RTX 4070} & \sNM; command in Section~\ref{sec:platform}\\\bottomrule
\end{tabular}
\end{table}

\paragraph{Historical performance (v2, not v4).} The only measured frame times in the repository are from the v2 renderer on an AMD Radeon integrated GPU at $1280\times720$
(\code{build/benchmark.json}, 29~Sep, 805 parts): v1 look 32.2\,FPS (31.0\,ms), LOW 18.2\,FPS, HIGH 14.7\,FPS (67.9\,ms), ULTRA 4.9\,FPS (202.7\,ms), shadow maps costing 16.2\,ms and 4$\times$ MSAA 9.3\,ms.
They describe a scene of 805 parts and a different renderer; they are given only so that the old report's table is not lost. \textbf{No v4 frame time exists}, neither on the integrated GPU nor on an RTX 4070, and the
EPIC preset has been exercised only on llvmpipe. The GPU timer labels in Section~\ref{sec:pipeline} are what \cli{benchmark} reports on the RTX machine. The part counts grew from 805 to \nLeaves{}
(the grass alone was 35\,000 instanced blades on HIGH, 450\,000 on EPIC), so the v2 numbers must not be extrapolated.

\subsection{Figures}
\begin{figure}[h]
  \centering
  \begin{subfigure}{0.495\linewidth}\picfour{base_gate}{\linewidth}\caption{Base gate}\end{subfigure}
  \begin{subfigure}{0.495\linewidth}\picfour{helipad}{\linewidth}\caption{Helipad}\end{subfigure}\\[1mm]
  \begin{subfigure}{0.495\linewidth}\picthree{cine_establishing}{\linewidth}\caption{v3 establishing shot of the compound}\end{subfigure}
  \begin{subfigure}{0.495\linewidth}\picthree{cine_chase}{\linewidth}\caption{v3 searchlight shaft (volumetric)}\end{subfigure}
  \caption{Results gallery: v4 base and the v3 cinematic shots (EPIC captures made on llvmpipe).}
  \label{fig:results}
\end{figure}

\begin{figure}[h]
  \centering
  \begin{subfigure}{0.325\linewidth}\pictwo{scene_xray}{\linewidth}\caption{X-ray edges (by primitive)}\end{subfigure}
  \begin{subfigure}{0.325\linewidth}\pictwo{scene_depth}{\linewidth}\caption{Z-buffer contents}\end{subfigure}
  \begin{subfigure}{0.325\linewidth}\pictwo{hide_ztest_off}{\linewidth}\caption{Depth test off (key \texttt{Z})}\end{subfigure}
  \caption{Diagnostic rendering (v2 captures). The Z-buffer view shows $z_\text{lin}=2nf/(f+n-z_\text{ndc}(f-n))$; hidden-line draws depth-tested edges over solids pushed back with
  \code{glPolygonOffset}, x-ray disables the depth test and colours all edges by primitive. The depth test keeps the rover hidden behind occluders (proposal \S3).}
  \label{fig:diag}
\end{figure}

\begin{figure}[h]
  \centering
  \begin{subfigure}{0.325\linewidth}\pictwo{compare_v1}{\linewidth}\caption{v1 renderer (key \texttt{N})}\end{subfigure}
  \begin{subfigure}{0.325\linewidth}\pictwo{compare_v2}{\linewidth}\caption{full pipeline}\end{subfigure}
  \begin{subfigure}{0.325\linewidth}\pictwo{ablate_shadows}{\linewidth}\caption{shadow maps off}\end{subfigure}\\[1mm]
  \begin{subfigure}{0.325\linewidth}\pictwo{scene_thermal_whitehot}{\linewidth}\caption{white-hot palette}\end{subfigure}
  \begin{subfigure}{0.325\linewidth}\pictwo{thermal_footprints}{\linewidth}\caption{warm tyre tracks}\end{subfigure}
  \begin{subfigure}{0.325\linewidth}\pictwo{rain_thermal}{\linewidth}\caption{thermal view in rain}\end{subfigure}
  \caption{v2 carry-over: same camera and time with the v1 look, the full pipeline and shadows off (surface bunker); thermal palettes, footprints and rain-cooled ground.}
  \label{fig:compare}
\end{figure}

\subsection{Limitations}
\begin{itemize}
  \item \textbf{Verification environment.} Every capture and test run was produced with Mesa llvmpipe (a CPU renderer); no run on a real GPU, no frame-time measurement, no NVENC encoding test.
  The visual tuning (explosion brightness, interior exposure) was done on llvmpipe images.
  \item \textbf{Windows.} The Windows executable is cross-compiled with zig; it has never been run natively, and the committed executable predates v4.
  \item \textbf{Soft particles.} The effect shader supports a soft depth fade but the renderer passes \code{uHasSceneDepth=0}, so smoke clips hard against geometry.
  \item \textbf{Combat scope.} Hit tests use spheres (the rovers, radius 0.4\,m) and boxes (static solids), not the part meshes. The tank shell and the AA tracers use the same
  target spheres as the UAV missile, but the player must aim them by hand; the player's machines have no damage state other than camouflage scorching, and nothing can hit the UAV.
  \item \textbf{Behaviour statistics.} The tracking statistics of v2 (TRACK 57\,\%, detection 51\,\%) were not re-measured with three intruders that are destroyed and respawn; \code{-{}-trace} is available.
  \item \textbf{Exporter.} The exporter labels the new surfaces 12--18 as ``Other'' and the terrain mesh function string with radius 20 (stale labels), and the workbook is updated by another engineer.
\end{itemize}

\subsection{Conclusion}
ThermalSense v4 turns the v3 patrol compound into a night military base without leaving the lab's rules: \nLeaves{} parts, all unit cubes, spheres, cylinders and cones with $T\,R\,S$ transforms,
share six indexed meshes. It adds an underground, walkable bunker with eased blast doors, an apron with a HESCO wall and an animated flag, windsock and gate, and six more machines with
non-trivial motion (arc-length track links, skid steer with a stabilised turret, a cyclic rotor with a spring-damper attitude, a spooling Gatling gun, three radar modes). The UAV's thermal vision drives
a complete engagement chain (lock, launch rules, proportional-navigation missile, blast, rigid-body debris, burning wreck, respawn) that is deterministic and covered by \nTests{} CPU tests, all passing.
The remaining work is measurement, not construction: FPS on the RTX 4070, a native Windows run, and a re-measured behaviour trace. The audit in Appendix~\ref{app:audit} lists every item and its evidence.
